% Generated draft: requires mathematical and source review.
\documentclass[11pt]{article}
\usepackage[margin=25mm]{geometry}
\usepackage{fontspec}
\setmainfont{Latin Modern Roman}
\usepackage{amsmath,amssymb,mathtools,unicode-math}
\setmathfont{Latin Modern Math}
\usepackage{enumitem,needspace,xurl,hyperref}
\setlength{\parindent}{0pt}
\setlength{\parskip}{5pt}
\setlength{\emergencystretch}{3em}
\newcommand{\N}{\mathbb{N}}
\newcommand{\Z}{\mathbb{Z}}
\newcommand{\Q}{\mathbb{Q}}
\newcommand{\R}{\mathbb{R}}
\newcommand{\C}{\mathbb{C}}
\newcommand{\F}{\mathbb{F}}
\newcommand{\A}{\mathbb{A}}
\newcommand{\PP}{\mathbb{P}}
\newcommand{\E}{\mathbb{E}}
\newcommand{\eps}{\varepsilon}
\newcommand{\dd}{\,\mathrm{d}}
\begin{document}
\section*{UnsolvedMath ID: 30005779}
\textbf{Source problem code:} catalogue entry
\subsection{Definitions and mathematical statement}
The supplied record is motivational and is marked \texttt{statement\_status: unrecoverable}. Its verification note says that the broad spectrum question motivates three questions preserved in a following record. Thus the material below records the available mathematical framework and questions, with unresolved hypotheses made explicit; it does not recover a uniquely specified standalone conjecture.

The supplied literature assessment suggests the setting of a finite-dimensional Hopf algebra \(A\) over a field \(k\). This is a contextual interpretation: the quoted passage does not specify the class of algebras, the field, its characteristic, or any additional tensor hypotheses. In this setting, let \(\operatorname{mod}_{\mathrm{fd}}A\) denote the category of finite-dimensional left \(A\)-modules, and define the stable category by
\[
\operatorname{stab}A
 =\operatorname{mod}_{\mathrm{fd}}A/[\operatorname{proj}A],
\]
where morphisms are taken modulo those factoring through projective modules. Finite-dimensional Hopf algebras are self-injective, so projective and injective modules coincide and this stable category is triangulated. The Hopf structure supplies the tensor product over \(k\), with unit the trivial module \(k\). Set
\[
H^*(A,k)=\bigoplus_{n\geq 0}\operatorname{Ext}^n_A(k,k),
\]
with its cohomological grading and Yoneda product.

For an essentially small symmetric tensor triangulated category \(\mathcal T\), the usual Balmer spectrum consists of its proper prime thick tensor ideals:
\[
\operatorname{Spc}(\mathcal T)
 =\{\mathcal P\subsetneq\mathcal T:\mathcal P\text{ is a thick tensor ideal and }
 X\otimes Y\in\mathcal P\Rightarrow X\in\mathcal P\text{ or }Y\in\mathcal P\}.
\]
Here thick means triangulated and closed under direct summands, and a tensor ideal is closed under tensoring with arbitrary objects. Its topology has closed subsets generated by the object supports
\[
\operatorname{supp}_{\mathcal T}(X)
 =\{\mathcal P\in\operatorname{Spc}(\mathcal T):X\notin\mathcal P\}.
\]
A general finite-dimensional Hopf algebra need not supply a symmetric tensor category. Consequently, this definition cannot be imposed on the entire contextual class without an additional hypothesis or a specified alternative spectrum. The source writes \(\operatorname{Spec}(\operatorname{stab}A)\); its intended spectrum convention remains unresolved.

The source displays a comparison map
\[
\Psi:\operatorname{Spec}(\operatorname{stab}A)
 \longrightarrow\operatorname{Proj}H^*(A,k).
\]
Under the usual graded-commutative interpretation, the underlying points of \(\operatorname{Proj}H^*(A,k)\) are homogeneous prime ideals not containing the positive-degree ideal, with the usual Zariski topology. The passage does not specify the construction of \(\Psi\), its complete hypotheses, or the convention under which this target is a scheme.

\emph{Source question.} For the algebras \(A\) intended by the report, how can one compute \(\operatorname{Spec}(\operatorname{stab}A)\)? The passage develops this motivation into the following questions:
\begin{enumerate}
\item When is \(\operatorname{Proj}H^*(A,k)\) a scheme of finite type? The source adds the parenthetical question: when is \(H^*(A,k)\) Noetherian?
\item When is the displayed comparison map \(\Psi\) a homeomorphism?
\item Can one calculate \(\operatorname{Proj}H^*(A,k)\)?
\end{enumerate}
These are questions about identifying hypotheses and carrying out computations, rather than a supplied universal assertion. The finite-type base is not stated; finite type of \(\operatorname{Proj}H^*(A,k)\) and Noetherianity of the graded ring are retained as separate questions, without asserting their equivalence.

The passage asserts that \(\Psi\) is surjective when \(H^*(A,k)\) is Noetherian. This is preserved as a source assertion with unspecified ambient hypotheses, not as an independently verified theorem for every finite-dimensional Hopf algebra. The passage also invokes results for finite groups and finite group schemes over a field, and says that the second question was touched upon in N.~Andruskiewitsch's talk.
\subsection{Short English statement}
The record asks how to compute the tensor triangular spectrum of a stable representation category, using cohomology and a proposed comparison map. It motivates questions about cohomological finiteness, when that map is a homeomorphism, and explicit computation of the cohomological projective spectrum. The record is marked unrecoverable: the intended algebra class, spectrum convention, and comparison-map hypotheses require human review.
\subsection{Sources}
\begin{itemize}
\item Supplied local UnsolvedMath registry record, permanent ID \(30005779\), title \emph{Balmer-Spectrum Motivation}, exact source problem code \texttt{catalogue entry}, proposed year 2024. No proposer is supplied. The record carries the supplied label \texttt{open} and the qualification \texttt{statement\_status: unrecoverable}; its verification note identifies this passage as motivation for three questions preserved in a following record.
\item Supplied report citation: \emph{Mini-Workshop: Bridging Number Theory and Nichols Algebras via Deformations} (2024), Oberwolfach Reports, pp.~235--272, DOI 10.4171/owr/2024/5. Source URL and literature-source link labelled ``original OWR report'': \url{https://doi.org/10.4171/owr/2024/5}. This authoritative report reference is supplied; it has not been independently consulted for this draft.
\item Supplied literature-source link labelled ``Nakano--Vashaw--Yakimov, spectrum and support theory'': \url{https://arxiv.org/abs/2112.11170}. The record supplies no further bibliographic details or exact result connecting this reference to the displayed map.
\item The quoted passage attributes tensor triangular geometry to Balmer in 2005, using its internal reference [1], and invokes finite-group and finite-group-scheme results through internal references [4] and [7]. Their bibliographic entries are not supplied. It also mentions N.~Andruskiewitsch's talk in connection with the homeomorphism question, without identifying the speaker of the quoted passage.
\item Literature assessment dated exactly \texttt{2026-08-22}: the supplied curation labels the general finite-dimensional Hopf-algebra spectrum question open and reports comparison maps and computations under finiteness hypotheses or for special tensor categories. This is supplied dated curation, not independently verified literature assessment. The background expressly requests independent verification of time-sensitive status and cited claims.
\end{itemize}
\end{document}
